\documentclass[12pt,a4paper]{ctexart}
\usepackage{amsmath,amssymb}
\usepackage{booktabs}
\usepackage{array}
\usepackage{geometry}
\usepackage{enumitem}
\usepackage{graphicx}
\usepackage[hidelinks]{hyperref}
\usepackage{fancyhdr}
\geometry{margin=2.5cm}
\setlength{\parskip}{0.5em}
\setlength{\parindent}{2em}
\linespread{1.25}
\renewcommand{\arraystretch}{1.15}


\pagestyle{fancy}
\fancyhf{}
\fancyfoot[C]{\normalfont\thepage}
\renewcommand{\headrulewidth}{0pt}
\renewcommand{\footrulewidth}{0pt}
\fancypagestyle{plain}{
  \fancyhf{}
  \fancyfoot[C]{\normalfont\thepage}
  \renewcommand{\headrulewidth}{0pt}
  \renewcommand{\footrulewidth}{0pt}
}

% ——中文论文格式规范:标题层级用黑体,正文宋体小四,表格五号——
\setcounter{secnumdepth}{-1}
\ctexset{
  section = {format = {\heiti \zihao{3}}},
  subsection = {format = {\heiti \zihao{4}}},
  subsubsection = {format = {\heiti \zihao{-4}}},
}

\newcommand{\dd}{\,\mathrm{d}}
\newcommand{\pdd}[2]{\dfrac{\partial #1}{\partial #2}}

\begin{document}

\begin{center}
{\heiti\zihao{-2} 2026 年高教社杯全国大学生数学建模竞赛\par}
\vspace{0.5em}
{\heiti\zihao{3} A 题\quad 药材的烘干问题\par}
\end{center}

\vspace{0.6em}
\begin{center}
{\heiti\zihao{3} 摘\quad 要}
\end{center}

本文针对中药材热风干燥的四个问题，建立了圆柱轴对称热质耦合传递模型，并给出离散通量严格闭合的数值求解方法。

\textbf{针对问题 1}，采用附录 2 的常物性参数与经验扩散系数 $D=7\times10^{-9}\mathrm{e}^{-0.89/C}$，建立圆柱径向一维热传导方程与水分扩散方程，表面取第三类（Robin）边界条件，烘房温度与水分浓度由附件 1 时变数据线性插值给出。用守恒型有限体积离散与 Crank--Nicolson 格式求解，得 1800 s 时表面温度由初值 28 $^\circ$C 升至 36.79 $^\circ$C，表面水分浓度由 2.55 kg/kg 降至 1.51 kg/kg，水分变化集中于近表面约 5 mm 的边界层内，中心含水率几乎不变，符合预热平衡阶段的物理特征。

\textbf{针对问题 2}，采用附录 3 的经验公式，其中密度、比热容、导热系数随水分浓度 $C$ 变化，扩散系数 $D(C,T)=2.4\times10^{-3}\mathrm{e}^{-0.45/C}\mathrm{e}^{-3850/T}$ 随 $C$ 与绝对温度 $T$ 变化，建立热质双向弱耦合模型：温度场通过 Arrhenius 因子影响水分扩散，水分场通过物性参数影响温度场；求解 0--3 h 的温度与水分场（表 3、表 4）。\textbf{针对问题 3}，沿用该模型，烘房条件取附件 1 插值（$t\le14400$ s）与恒温段参数 $(50.165\,^\circ\mathrm{C},\,0.04986\ \mathrm{kg/kg})$ 两段拼接，按判据「各处水分浓度 $\le0.15$ kg/kg」求得烘干时长 $t^*=205701$ s，即 \textbf{57.14 h}（约 2.38 天），处于题述「烘干过程一般持续 2--3 天」范围内。

\textbf{针对问题 4}，考虑药材尺寸收缩：采用附录 4 的经验公式，并以附件 2 的半径数据经 Pchip 保形插值构造动边界 $R(t)$；引入随体坐标 $\zeta=r/R(t)$，将热传导与水分扩散方程变换为无对流项的 Lagrange 形式，在同一守恒离散框架下求解。求得含收缩时的烘干时长 $t^*=182897$ s，即 \textbf{50.80 h}（约 2.12 天）。

\textbf{数值方法与检验}：空间采用离散通量严格闭合的有限体积离散（表面节点位于真实表面，Robin 通量与内部界面通量逐时间层对消），时间采用 Crank--Nicolson 格式，非线性与双向耦合用交错 Picard 迭代处理。模型通过四重检验：与 Bessel 级数解析解对照，温度场四位小数完全一致；空间步长减半收敛比达理论值 4（问题 1、2 实测比值 4.0），验证二阶精度；离散质量守恒闭合差在 $10^{-7}$ 以内（问题 1 达机器精度）；长时程渐近精确收敛到烘房参数。此外对水分汽化潜热作数量级分析：表面蒸发吸热与对流显热之比约 3--5，其不确定性以时长敏感性区间（约 $+40\%\sim50\%$，见 7.6 节）计入结果解读。

\medskip
\noindent{\heiti 关键词}：热风干燥；热质耦合传递；圆柱轴对称扩散；Crank--Nicolson 格式；移动边界；有效水分扩散系数

\newpage
\section{一、问题重述}

中药材热风干燥通常包含预热平衡与恒温干燥两个阶段，通过调控烘房温湿环境完成干燥。某中药材形状近似为圆柱形，长 25 cm，半径 2 cm，初始温度 28 $^\circ$C，初始水分浓度（干基含水率）2.55 kg/kg。需建立数学模型研究如下四个问题：

\textbf{问题 1}：附件 1 给出烘干初期烘房温度与水分浓度的时变数据，附录 2 给出相关参数。建立预热平衡阶段药材温度与水分浓度变化规律的数学模型，按表 1、表 2 格式给出 100--1800 s 内 7 个时刻、距中心 0--2 cm 处的结果，并给出 1800 s 内每隔 1 s、距离每隔 0.1 cm 的完整结果。

\textbf{问题 2}：整个烘干过程包含参数不同的预热平衡与恒温干燥两阶段，经验公式统一采用附录 3。建立整个烘干过程的数学模型，按表 3、表 4 格式给出前 3 h 内每隔 0.5 h、距中心 0--2 cm 处的结果，并给出每隔 1 s、距离每隔 0.1 cm 的完整结果。

\textbf{问题 3}：按烘干要求，药材各处水分浓度应低于 0.15 kg/kg。确定烘干所需时间，按表 5 格式给出每隔 6 h、距中心每隔 0.5 cm 处的水分浓度，并给出每隔 60 s、距离每隔 0.1 cm 的完整结果。

\textbf{问题 4}：实际烘干中药材因水分流失发生尺寸变化。根据附件 2 的半径数据与附录 4 的经验公式，确定烘干时长，按表 6 格式给出每隔 6 h、距中心每隔 0.5 cm 处及药材表面处的水分浓度，并给出每隔 60 s、距离每隔 0.1 cm 的完整结果。

所有结果保留四位小数。

\section{二、问题分析}

\subsection{2.1 几何与传递机理分析}

药材为长 25 cm、半径 2 cm 的细长圆柱，长径比达 6.25。轴向与径向的扩散特征时间之比约为 $(L/(2R))^2\approx39$（轴向特征尺度取半长 $L/2$，即到最近端面的距离），故端面效应与轴向不均匀性可以忽略，药材内部的温度与水分梯度主要由径向热风对流引起。将药材视为无限长圆柱，取径向一维轴对称模型即可（文献中圆柱形食品与药材的干燥均广泛采用径向一维模型\cite{simal,abbasi}）。周向按轴对称处理，初始条件均匀，故解不依赖周向角。

内部传递机理：热量按傅里叶定律自表面向内传导\cite{incropera,yang}；水分在浓度梯度驱动下按菲克定律向表面扩散，在表面以对流方式传递到烘房空气。表面边界均为第三类（Robin）条件：热侧由对流换热系数 $h$ 描述，质侧由对流传质系数 $h_m$ 描述。

\subsection{2.2 热质耦合的机制分析}

本题的耦合是\textbf{双向弱耦合}：温度场影响水分场——扩散系数经验式 $D(C,T)$ 含 Arrhenius 因子 $\mathrm{e}^{-3850/T}$，温度升高显著加速水分扩散；水分场影响温度场——附录 3、4 中密度、比热容、导热系数均随水分浓度 $C$ 变化，水分流失改变热容与热阻。两个场通过物性参数相互影响，而边界条件本身不直接交叉（无 Soret/Dufour 交叉效应）。数值上采用交错 Picard 迭代：同一时间步内先由当前水分场求温度场，再用新温度场求水分场，交错循环直至满足 5.5 节的收敛判据（各算例平均 2.2--3.4 轮，最多 4 轮）。

关于水分汽化的局部吸热效应：严格意义下蒸发吸热会降低表面温度，但本题题给的经验式（$D(C,T)$、$h_m$ 等）是由干燥实验数据拟合得到的\textbf{有效参数}，其数值可能已隐含吸收汽化与吸附等微观机制的净效应——现有信息不足以唯一判定其吸收程度，但若在题给经验式之上再显式加入潜热汇项，将面临与经验参数重复计账的风险，且需引入题给之外的饱和蒸气压等参数。这一判断的数量级依据在 7.6 节给出专门分析（时长敏感性区间约 $+40\%\sim50\%$），并在模型假设 H5 中予以声明。

\subsection{2.3 两阶段烘房条件的处理}

附件 1 给出烘干初期（0--14400 s）烘房温度与水分浓度的时变序列（共 241 个数据点）。整个烘干过程分为预热平衡与恒温干燥两阶段：预热平衡段烘房条件按附件 1 线性插值；其后进入恒温段，烘房参数取附件 1 末值并保持恒定，即 $T_a=50.165\ ^\circ$C，$C_a=0.04986$ kg/kg。问题 1 的时间范围（0--1800 s）与问题 2 的输出范围（前 3 h）均落在附件 1 数据段内，直接插值即可；问题 3、4 的求解需贯穿全程，采用上述两段拼接的边界条件。

\subsection{2.4 四个问题的建模差异}

四个问题在物性经验式、烘房条件、几何处理与输出要求上的差异是本题的关键，逐一对照如下：

\begin{center}
{\zihao{5}\begin{tabular}{lcccc}
\toprule
条件 & 问题 1 & 问题 2 & 问题 3 & 问题 4 \\
\midrule
物性经验式 & 附录 2 & 附录 3 & 附录 3 & 附录 4 \\
$D$ 是否随 $T$ 变 & 否（仅随 $C$） & 是 & 是 & 是 \\
$\rho,c_p,k$ 是否随 $C$ 变 & 否（常数） & 是 & 是 & 是 \\
烘房条件 & 附件 1 & 附件 1 & 附件 1$+$恒温段 & 附件 1$+$恒温段 \\
几何收缩（附件 2） & 否 & 否 & 否 & 是（$R(t)$） \\
求解时长 & 1800 s & 前 3 h & 至 $C\le0.15$ & 至 $C\le0.15$ \\
输出 & 表 1/2 & 表 3/4 & 表 5、$t^*$ & 表 6、$t^*$ \\
\bottomrule
\end{tabular}}
\end{center}
\section{三、模型假设}

\begin{enumerate}[label=H\arabic*.,leftmargin=2.2em]
\item \textbf{连续介质近似。}药材内部的多孔骨架、细胞间隙与水分经体积平均视为连续介质，温度与水分浓度在宏观尺度上连续分布，可用傅里叶定律与菲克定律描述。这一做法是多孔介质干燥理论的标准框架\cite{whitaker,mujumdar}。
\item \textbf{径向一维轴对称。}长径比 6.25，忽略轴向传递与端面效应，忽略周向不均匀性，温度与水分浓度仅为 $r$（或 $\zeta$）与 $t$ 的函数。
\item \textbf{初值均匀。}$T(r,0)=28\ ^\circ$C，$C(r,0)=2.55$ kg/kg。
\item \textbf{对流系数恒定。}$h=25$ W/(m$^2\!\cdot$K)、$h_m=8\times10^{-7}$ m/s 在全程视为常数。文献研究表明，常规对流干燥中对流系数在恒速段与降速段内近似为常数，仅在段间过渡处显著变化\cite{mujumdar,defraeye}，故该简化合理。
\item \textbf{经验式的有效参数观点（含汽化潜热）。}题给的 $\rho,c_p,k,D,h_m$ 是由实验标定的有效参数：$D(C,T)$ 的 Arrhenius 指数对应的活化能 $E_a=32.0$ kJ/mol 高于纯液态水自扩散的活化能（约 18--20 kJ/mol\cite{mills}），其差值对应解吸与汽化的综合能垒，说明潜热效应可能已被有效参数部分吸收，但现有信息不足以唯一判定其吸收程度（数量级分析见 7.6 节）。因此主模型采用显热传导与有效水分扩散系数，不显式引入蒸发潜热汇项：若显式叠加，将面临与经验参数重复计账的风险，且需引入题给之外的饱和蒸气压等数据。
\item \textbf{问题 1--3 忽略收缩。}题面仅要求问题 4 依据附件 2 考虑尺寸变化，且附录 2、3 的物性参数不含收缩机制，故问题 1--3 取固定半径 $R=2$ cm。附件 2 数据显示前 30 min 半径收缩约 6.4\%，全程收缩约 40.1\%，该简化对问题 1--3 的影响在问题 4 中通过含收缩模型予以评估（见 6.4 节对比）。
\item \textbf{烘房环境均匀且按题给数据变化。}预热平衡段按附件 1 线性插值；恒温干燥段取附件 1 末值 $(50.165\ ^\circ$C$, 0.04986\ \mathrm{kg/kg})$ 恒定——14400 s 后固定为附件 1 末值是对题给信息的重要外延（附件 1 仅覆盖 0--14400 s），其对结论的影响由 7.7 节的参数敏感性分析评估。
\item \textbf{收缩各向同性。}问题 4 中半径按附件 2 收缩，且径向收缩均匀（水分浓度与温度仅沿径向变化，与题给一维数据自洽）；干基含水率定义于固相骨架，随体坐标下无对流项（见 5.4 节）。
\end{enumerate}

\section{四、符号说明}

\begin{center}
{\zihao{5}\begin{tabular}{llll}
\toprule
符号 & 含义 & 符号 & 含义 \\
\midrule
$r$ & 径向坐标（m） & $t$ & 时间（s） \\
$R$ & 药材半径（m，问题 4 中为 $R(t)$） & $\zeta$ & 随体坐标 $r/R(t)$ \\
$T$ & 药材温度（$^\circ$C） & $T_a$ & 烘房温度（$^\circ$C） \\
$C$ & 水分浓度/干基含水率（kg/kg） & $C_a$ & 烘房水分浓度（kg/kg） \\
$\rho$ & 密度（kg/m$^3$） & $c_p$ & 比热容（J/(kg$\cdot$K)） \\
$k$ & 导热系数（W/(m$\cdot$K)） & $D$ & 水分扩散系数（m$^2$/s） \\
$h$ & 对流换热系数（W/(m$^2\!\cdot$K)） & $h_m$ & 对流传质系数（m/s） \\
$T_0,C_0$ & 初始温度/水分浓度 & $t^*$ & 烘干时长（h） \\
$\Delta r,\Delta t$ & 空间/时间步长 & $N$ & 径向节点数 \\
$\mathrm{Bi}$ & 毕渥数 $hR/k$ & $E_a$ & 表观活化能（kJ/mol） \\
$\Delta H_v$ & 水的汽化潜热（J/kg） & $\mathrm{RH}$ & 空气相对湿度 \\
\bottomrule
\end{tabular}}
\end{center}

\section{五、模型的建立与求解}

\subsection{5.1 问题 1：预热平衡阶段的模型}

预热平衡阶段药材尚处低温（初温 28 $^\circ$C），物性可按附录 2 取常数：$\rho=820$ kg/m$^3$，$c_p=2600$ J/(kg$\cdot$K)，$k=0.36$ W/(m$\cdot$K)，$h=25$ W/(m$^2\!\cdot$K)，$h_m=8\times10^{-7}$ m/s；水分扩散系数采用经验式
\begin{equation}
D(C)=7\times10^{-9}\,\mathrm{e}^{-0.89/C}\quad(\mathrm{m^2/s}).
\label{eq:D1}
\end{equation}
径向一维轴对称模型的控制方程、定解条件为
\begin{align}
\rho c_p\pdd{T}{t}&=\frac{k}{r}\pdd{}{r}\Bigl(r\pdd{T}{r}\Bigr), & 0<r<R,\ t>0, \label{eq:heat1}\\
\pdd{C}{t}&=\frac{1}{r}\pdd{}{r}\Bigl(r\,D(C)\pdd{C}{r}\Bigr), & 0<r<R,\ t>0, \label{eq:mass1}\\
\pdd{T}{r}(0,t)&=0,\qquad -k\pdd{T}{r}(R,t)=h\bigl(T(R,t)-T_a(t)\bigr), \label{eq:bcT1}\\
\pdd{C}{r}(0,t)&=0,\qquad -D\pdd{C}{r}(R,t)=h_m\bigl(C(R,t)-C_a(t)\bigr), \label{eq:bcC1}\\
T(r,0)&=T_0=28\ ^\circ\mathrm{C},\qquad C(r,0)=C_0=2.55\ \mathrm{kg/kg}. \label{eq:ic1}
\end{align}
其中 $T_a(t)$、$C_a(t)$ 由附件 1 线性插值；附件 1 的 $C_a$ 单位为 kg/kg、与干基含水率 $C$ 同纲，本文将其解释为烘房空气状态对应的药材等效平衡水分浓度，故 Robin 边界驱动力 $C_R-C_a$ 直接成立（按空气含湿比解释的情景分析见 7.6 节）。式 \eqref{eq:heat1}、\eqref{eq:mass1} 在 $r=0$ 处按对称性与洛必达法则化为 $\pdd{T}{t}=2\alpha\,\pdd{^2T}{r^2}$ 及 $\pdd{C}{t}=2\pdd{}{r}\!\bigl(D(C)\pdd{C}{r}\bigr)$，其中 $\alpha=k/(\rho c_p)$ 为热扩散率。

\subsection{5.2 问题 2：0--3 h 温度与水分场模型}

物性按题设统一采用附录 3 的经验式：
\begin{equation}
\rho(C)=650+128C,\quad c_p(C)=1450+2736\frac{C}{C+1},\quad k(C)=0.21+0.38\frac{C}{C+1},
\label{eq:props3}
\end{equation}
\begin{equation}
D(C,T)=2.4\times10^{-3}\,\mathrm{e}^{-0.45/C}\,\mathrm{e}^{-3850/T}\quad(\mathrm{m^2/s}),
\label{eq:D3}
\end{equation}
其中 $T$ 为绝对温度（K）。此时热传导方程中的 $\rho c_p$ 与 $k$ 均随 $C$ 变化，且 $D$ 同时依赖 $C$ 与 $T$，两场构成\textbf{双向弱耦合}（边界条件不直接交叉，且忽略水分输运携带的显焓等交叉项，见 2.2 节）：
\begin{align}
\rho(C)c_p(C)\pdd{T}{t}&=\frac{1}{r}\pdd{}{r}\Bigl(r\,k(C)\pdd{T}{r}\Bigr), \label{eq:heat23}\\
\pdd{C}{t}&=\frac{1}{r}\pdd{}{r}\Bigl(r\,D(C,T)\pdd{C}{r}\Bigr), \label{eq:mass23}
\end{align}
边界条件与式 \eqref{eq:bcT1}、\eqref{eq:bcC1} 形式相同，其中质侧条件的 $D$ 取表面值 $D(C_R,T_R)$。题面未另行给出对流系数，$h$、$h_m$ 沿用附录 2 取值（$h=25$ W/(m$^2\!\cdot$K)，$h_m=8\times10^{-7}$ m/s），与假设 H4 一致。初值同式 \eqref{eq:ic1}。

烘房条件 $T_a(t)$、$C_a(t)$ 按 2.3 节由附件 1 线性插值给出，两场耦合机制见 2.2 节，求解结果按表 3、表 4 输出。

\subsection{5.3 问题 3：全程烘干模型与时长判据}

问题 3 沿用式 \eqref{eq:heat23}--\eqref{eq:mass23} 的同一方程组与初边值（物性及 $h$、$h_m$ 同 5.2 节），但需求解整个烘干过程，时间范围远超附件 1 的 14400 s，故烘房条件按两阶段拼接：
\begin{equation}
T_a(t)=
\begin{cases}
\text{附件 1 线性插值}, & t\le14400\ \mathrm{s},\\
50.165\ ^\circ\mathrm{C}, & t>14400\ \mathrm{s},
\end{cases}
\quad
C_a(t)=
\begin{cases}
\text{附件 1 线性插值}, & t\le14400\ \mathrm{s},\\
0.04986\ \mathrm{kg/kg}, & t>14400\ \mathrm{s}.
\end{cases}
\label{eq:twostage}
\end{equation}
即前 14400 s 按附件 1 线性插值，其后进入恒温段，烘房参数取附件 1 末值并保持恒定。按题面判据求解：当「各处水分浓度 $\le0.15$ kg/kg」即 $\max_r C(r,t)\le0.15$ kg/kg 首次成立时停止，对应时刻即烘干时长 $t^*$，全程水分场按表 5 输出。

\subsection{5.4 问题 4：含收缩的模型}

\textbf{(1) 动边界与随体坐标。}附件 2 给出每隔 1800 s 的药材半径数据（共 145 点，覆盖 0--72 h，自 67 h 起恒定于 1.198 cm）。为得到光滑的 $R(t)$，采用保形分段三次 Hermite（Pchip）插值（保单调、无过冲），并取 $\mathrm{d}R/\mathrm{d}t$ 为其解析导数。收缩使外边界随物料移动，故将径向坐标换为随体（Lagrange）坐标
\begin{equation}
\zeta=\frac{r}{R(t)}\in[0,1],
\label{eq:zeta}
\end{equation}

各向同性且径向均匀的收缩（假设 H8）意味着任一物质点按统一比例 $R(t)/R(0)$ 收缩，即物质点位置 $r_p(t)=\zeta_pR(t)$、径向速度 $v_r=\zeta_pR'(t)=(R'/R)\,r_p$，故 $\zeta$ 恒标记同一干物质质点。下面从守恒律推导随体坐标下的方程：干基含水率 $C$ 定义于固相骨架，对随骨架运动的物质元，干固体质量守恒要求 $\mathrm{d}M_d=\rho_d\,\mathrm{d}V$ 恒定（$\rho_d=\rho/(1+C)$ 为干物质密度），水分守恒给出 $\mathrm{d}M_d\,\mathrm{D}C/\mathrm{D}t=-\nabla\cdot q_d\,\mathrm{d}V$，其中 $q_d=-\rho_d D\nabla C$ 为相对骨架的 Fick 通量\cite{whitaker}，$\mathrm{D}/\mathrm{D}t$ 为随体导数。联立两式并利用干固体守恒恒等式 $\partial\rho_d/\partial t+\nabla\cdot(\rho_d v)=0$，对流项 $v\cdot\nabla C$ 与 $\rho_d$ 的变化率项严格相消，得 $\rho_d\,\mathrm{D}C/\mathrm{D}t=\nabla\cdot(\rho_d D\nabla C)$——这正是「随体坐标无对流项」的由来。在题给有效参数观点下（假设 H5），$\rho_d$ 随 $C$ 的变化已隐含于 $D(C,T)$ 的经验标定中，取 $\rho_d$ 均匀近似即得式 \eqref{eq:mass4} 的水分方程；热量方程经同样的随体变换并忽略水分携焓交叉项（见 2.2 节）化为
\begin{equation}
\pdd{C}{t}\Big|_{\zeta}=\frac{1}{R^2(t)\,\zeta}\pdd{}{\zeta}\Bigl(\zeta\,D(C,T)\pdd{C}{\zeta}\Bigr),
\label{eq:mass4}
\end{equation}
\begin{equation}
\rho(C)c_p(C)\pdd{T}{t}\Big|_{\zeta}=\frac{1}{R^2(t)\,\zeta}\pdd{}{\zeta}\Bigl(\zeta\,k(C)\pdd{T}{\zeta}\Bigr),
\label{eq:heat4}
\end{equation}

物性由附录 4 给出：$\rho=760+90C$，$c_p=1850+2150\,C/(C+1)$，$k=0.12+0.20\,C/(C+1)$，
\begin{equation}
D(C,T)=4.2\times10^{-4}\,\mathrm{e}^{-0.30/C}\,\mathrm{e}^{-3850/T}\quad(\mathrm{m^2/s}).
\label{eq:D4}
\end{equation}
边界条件（$\zeta=1$ 为真实表面）：
\begin{equation}
-\frac{k}{R}\pdd{T}{\zeta}(1,t)=h\bigl(T(1,t)-T_a(t)\bigr),\qquad
-\frac{D}{R}\pdd{C}{\zeta}(1,t)=h_m\bigl(C(1,t)-C_a(t)\bigr),
\label{eq:bc4}
\end{equation}
对称条件 $\pdd{T}{\zeta}(0,t)=\pdd{C}{\zeta}(0,t)=0$。表 6 中「到药材中心的距离」列在 $d>R(t)$ 时该位置已在药材之外，表中以 ``---'' 表示（相应单元格按模板留空）；「药材表面」列即 $\zeta=1$ 处的值。固定距离列在 $\zeta$ 坐标的离散节点上经 Pchip 插值取值（插值阶的选择依据见 7.2 节）。

\textbf{(2) 收缩对烘干的影响机制。}由式 \eqref{eq:mass4} 可见收缩的加速作用是同一几何效应的两种等价表述：$R(t)$ 减小使算子系数 $D/R^2$ 增大（全程 $R$ 由 2 cm 减至 1.198 cm，$R^{-2}$ 增大 2.8 倍），其物理含义即浓度场随骨架向内收缩、等效扩散路径缩短，两者并非相互独立的机制。附录 4 的 $D$ 前因子与指数均小于附录 3，收缩效应与物性差异的综合影响在 6.4 节对比。

\subsection{5.5 数值方法}

四个问题共用一套守恒型离散框架，下面以问题 1 的径向空间离散（问题 4 的 $\zeta$ 坐标离散同理）说明。

\textbf{(1) 空间离散节点与控制体。}取节点 $r_i=i\Delta r$（$i=0,1,\dots,N-1$），则表面节点恰位于真实表面：$r_{N-1}=R$。每个节点 $i$ 对应控制体体积 $v_i=\int_{CV_i} r\,\mathrm{d}r$：内部节点取 $[r_{i-1/2},r_{i+1/2}]$，圆心节点取 $[0,r_{1/2}]$，表面节点取 $[r_{N-3/2},R]$。方程 \eqref{eq:heat1}--\eqref{eq:mass1} 在控制体上积分（对 $r\dd r$），通量在界面 $i\pm1/2$ 用算术平均的扩散系数（如 $D_{i-1/2}=\frac12(D_{i-1}+D_i)$）离散，得到
\begin{equation}
\frac{u_i^{n+1}-u_i^n}{\Delta t}=\frac12\bigl(L u^{n+1}+Lu^n\bigr)_i+\frac12\bigl(s^{n+1}+s^n\bigr)_i,
\label{eq:cn}
\end{equation}
其中 $u$ 代表 $T$ 或 $C$，$(Lu)_i$ 为界面通量差
\begin{equation}
(Lu)_i=\Bigl[r_{i+1/2}K_{i+1/2}\frac{u_{i+1}-u_i}{\Delta r}-r_{i-1/2}K_{i-1/2}\frac{u_i-u_{i-1}}{\Delta r}\Bigr]/v_i,
\label{eq:L}
\end{equation}
圆心行按对称性与洛必达法则取 $4K_{1/2}(u_1-u_0)/\Delta r^2$（$v_0=\Delta r^2/8$）。表面控制体的体积为 $v_{N-1}=\bigl((N-1)^2-(N-\tfrac32)^2\bigr)\Delta r^2/2$，其外界面即真实表面，将 Robin 条件 \eqref{eq:bcT1}、\eqref{eq:bcC1} 直接作为外界面通量作用于表面节点：
\begin{equation}
(Lu)_{N-1}=\Bigl[r_{N-3/2}K_{N-3/2}\frac{u_{N-2}-u_{N-1}}{\Delta r}-R\,\beta\,K_s\,(u_{N-1}-u_a)\Bigr]/v_{N-1},
\label{eq:Lsurf}
\end{equation}
其中热方程 $\beta=h/k_s$，质方程 $\beta=h_m/D_s$。式 \eqref{eq:Lsurf} 中表面节点值即真实表面值，通量取 $R\,h\,(T_R-T_a)$ 与 $R\,h_m(C_R-C_a)$，与内部界面通量逐时间层严格对消，保证离散守恒；各离散项截断误差的泰勒展开分析见本节 (5)。

\textbf{(2) 时间推进与线性代数。}式 \eqref{eq:cn} 为 Crank--Nicolson 格式\cite{crank}（无条件稳定、时间二阶），每步化为三对角方程组，用追赶法（Thomas 算法）求解：对三对角方程组 $a_ix_{i-1}+b_ix_i+c_ix_{i+1}=d_i$，前代 $c'_0=c_0/b_0$、$d'_0=d_0/b_0$，随后 $m=b_i-a_ic'_{i-1}$、$c'_i=c_i/m$、$d'_i=(d_i-a_id'_{i-1})/m$；回代 $x_{N-1}=d'_{N-1}$、$x_i=d'_i-c'_ix_{i+1}$，复杂度 $O(N)$。空间通量项 $L$ 由式 \eqref{eq:L}、\eqref{eq:Lsurf} 给出，时间层上源项（边界贡献）取 $n$、$n+1$ 两层的算术平均。

\textbf{(3) 非线性与耦合的处理。}附录 3、4 的物性随 $C,T$ 变化，且两场双向耦合，采用\textbf{交错 Picard 迭代}：时间步 $[t_n,t_{n+1}]$ 内，
\begin{enumerate}[label=(\roman*),leftmargin=3em]
\item 用当前水分场 $C^{k}$ 计算 $k(C)$、$\rho c_p$，求解温度场得 $T^{k+1}$；
\item 用 $C^{k}$ 与温度 $T^{k+1}$ 计算 $D(C,T)$，求解水分场得 $C^{k+1}$；
\item 以 $k+1\to k$ 重复 (i)(ii)，直至收敛判据满足（系数在迭代中取最新可用层，时间平均仍按 Crank--Nicolson 取 $n$、$n+1$ 两层）。
\end{enumerate}
收敛判据取 $T$、$C$ 的无穷范数相对变化（按 $\|u^{k+1}-u^k\|_\infty/\max(\|u^{k+1}\|_\infty,1)$ 计）均小于 $10^{-8}$，最多迭代 20 轮。停止阈值的选取依据：$\Delta t$ 减半实测的时间步误差为 $O(10^{-5})$（7.2 节），迭代残差须再低 2--3 个量级方不成为精度瓶颈，故取 $10^{-8}$（相对量）。实测问题 2--4 全程平均 2.2--3.4 轮（最多 4 轮）即收敛，收敛轮数随步长与干燥阶段变化小，迭代误差不进入精度分析。同一格式经变量代换直接适用于问题 4 的 $\zeta$ 坐标：控制体体积改为 $\int\zeta\dd\zeta$，算子系数乘 $1/R^2(t)$，表面通量取 $h\,R\,(T_R-T_a)$ 与 $h_mR\,(C_R-C_a)$，$R$ 取半步中点值 $R(t_{n+1/2})$。

\textbf{(4) 离散守恒恒等式。}由式 \eqref{eq:cn}、\eqref{eq:L}、\eqref{eq:Lsurf} 对全体控制体求和，内部界面通量逐项对消（伸缩和），得逐时间层的严格恒等式
\begin{equation}
\sum_{i=0}^{N-1}v_i\bigl(C_i^{n+1}-C_i^n\bigr)
=-\frac{\Delta t}{2}R\Bigl[h_m\bigl(C_R^{n+1}-C_a^{n+1}\bigr)+h_m\bigl(C_R^n-C_a^n\bigr)\Bigr],
\label{eq:conserv}
\end{equation}
即域内水分增量恒等于表面传递量，该恒等式被用于 7.3 节的守恒检验，闭合差在 $10^{-7}$ 以内。问题 4 的 $\zeta$ 坐标离散同理：控制体体积取 $\int\zeta\dd\zeta$，界面通量系数为 $\zeta_{i\pm1/2}D_m/\Delta\zeta$，表面通量为 $h_mR(C_R-C_a)$。

\textbf{(5) 截断误差与二阶精度。}设解充分光滑（步长足够小时成立，数值证据见 7.2 节），各离散项的截断误差由泰勒展开逐项给出。

\textbf{(i) 空间通量项。}界面差分在 $r_{i+1/2}$ 处展开
\begin{equation}
\frac{u_{i+1}-u_i}{\Delta r}
=u'\Big|_{i+1/2}+\frac{\Delta r^2}{24}u'''+O(\Delta r^4),
\end{equation}
故界面通量 $r_{i+1/2}K_{i+1/2}(u_{i+1}-u_i)/\Delta r$ 的误差为 $O(\Delta r^2)$。相邻界面通量相减，主部为 $\Delta r\,(rKu')'$、误差项 $O(\Delta r^3)$；除以控制体体积 $v_i=r_i\Delta r$（内部节点）后，$(Lu)_i$ 的截断误差为 $O(\Delta r^2)$。圆心行由对称性 $u_1=u_0+\frac{\Delta r^2}{2}u''_0+O(\Delta r^4)$，得 $4K_{1/2}(u_1-u_0)/\Delta r^2=2Ku''_0+O(\Delta r^2)$，恰为 $(K/r)\partial_r(r\,\partial_r u)|_{r=0}=2Ku''$ 的二阶逼近；表面行外界面 Robin 通量直接采用表面节点真值，无数值微分误差，内界面仍为中心差分，故表面行同为二阶。

\textbf{(ii) 界面扩散系数的算术平均。}对光滑的 $D(r)$（由 $D(C(r),T(r))$ 与解的光滑性传递）：
\begin{equation}
\frac{D_i+D_{i+1}}{2}=D\Big|_{i+1/2}+\frac{\Delta r^2}{8}D''+O(\Delta r^4),
\end{equation}
该系数偏差对通量的贡献为 $O(\Delta r^2)$ 且随 $r$ 光滑变化，相邻界面相减时同样缩并，$(Lu)_i$ 的总截断误差仍为 $O(\Delta r^2)$。

\textbf{(iii) 时间项。}在 $t_{n+1/2}$ 处展开：
\begin{equation}
\frac{u^{n+1}-u^n}{\Delta t}=\pdd{u}{t}\Big|_{n+1/2}+O(\Delta t^2),\qquad
\frac12\bigl(Lu^{n+1}+Lu^n\bigr)=Lu\Big|_{n+1/2}+O(\Delta t^2),
\end{equation}
故单步截断误差为 $O(\Delta t^2+\Delta r^2)$，格式空间、时间均二阶。问题 4 的 $\zeta$ 坐标形式经变量代换，误差结构相同（算子系数多乘 $1/R^2(t)$）；动边界下 $R$ 取半步中点值额外引入 $O(\Delta t^2)$ 配对误差（见 7.3 节）。综上，格式的整体截断误差为 $O(\Delta r^2+\Delta t^2)$，与 7.2 节步长减半收敛比（问题 1、2 实测 4.0；陡梯度输出行收敛阶偏低，按实测阶估计）互相印证，支撑 5.6 节以误差估计选取步长的流程。

\subsection{5.6 空间离散与步长的选取}

\textbf{选取原则。}题面要求所有结果保留四位小数，故数值误差须低于四位小数末位的一半，即 $0.5\times10^{-4}$（否则末位不可信）。本文格式空间、时间均为二阶收敛（误差随步长二次下降），按如下流程选取步长：对每问先取一套初始步长求解，再与步长减半的解逐格比较，用两套解的差估计细步长解的误差（Richardson 估计，二阶收敛时 $E(h/2)\approx|u_h-u_{h/2}|/3$）；若估计值仍高于 $0.5\times10^{-4}$，则继续减半加密，直至降至 $10^{-5}$ 量级（较阈值留约一个量级的裕度），方确定为生产步长。需要说明的是，对表 5、表 6 的固定时刻行与 $t^*$ 等输出量，实测收敛阶受陡梯度区与插值取值影响而低于 2，此时 Richardson 估计改用实测阶（见问题 3、4）。

时间步长的选取由精度而非稳定性决定（Crank--Nicolson 格式无条件稳定）：要求时间离散误差比空间离散误差再小 1--2 个量级，不成为精度瓶颈，同样以 $\Delta t$ 减半的实测差值验证。

\textbf{问题 1。}先作物理尺度估计：表面浓度剖面的特征长度约 $\delta\sim D/h_m\approx5$ mm（$D$ 取表面典型值 $\sim4\times10^{-9}$ m$^2$/s，$h_m=8\times10^{-7}$ m/s，即 Robin 边界的外推长度），浓度梯度最陡处位于 $(100\ \mathrm{s},$ 表面$)$，特征长度内的指数型浓度分布需要足够多的节点才能分辨，故从 $\Delta r=0.125$ mm 起步。加密实测：$0.125\to0.0625$ mm 时表 2 最大差 $3.51\times10^{-4}$ kg/kg，即 0.0625 mm 解的 Richardson 估计约 $1.2\times10^{-4}$，仍高于 $0.5\times10^{-4}$；继续加密 $0.0625\to0.03125$ mm，最大差 $8.72\times10^{-5}$，0.03125 mm 解的估计约 $2.9\times10^{-5}$，已满足要求（较阈值留约半个量级裕度）。故问题 1 取 $\Delta r=0.03125$ mm（径向节点 $N=641$），$\Delta t=0.125$ s（$\Delta t$ 减半后表 2 差值为 $O(10^{-5})$，见 7.2 节，时间误差比空间误差小约一个量级）。

\textbf{问题 2。}输出自 0.5 h 开始，初始边界层瞬态已过，浓度分布较平缓，按同一流程加密 $0.5\to0.25\to0.125$ mm：表 4 最大差依次 $1.86\times10^{-4}$、$4.66\times10^{-5}$ kg/kg，0.125 mm 解的 Richardson 估计约 $1.6\times10^{-5}$，低于 $0.5\times10^{-4}$ 阈值（0.25 mm 解约 $6.2\times10^{-5}$，仍不达标），故取 $\Delta r=0.125$ mm（$N=161$），$\Delta t=0.25$ s（减半后表 3 最大差 $2.02\times10^{-5}\ ^\circ$C、表 4 最大差 $6.59\times10^{-6}$ kg/kg）。

\textbf{问题 3。}需求解至约 57 h，且烘干时长 $t^*$ 由末行判据决定、对空间离散更敏感，故在问题 2 基础上继续加密。实测 $0.5\to0.25\to0.125$ mm：表 5 最大差依次 $5.27\times10^{-4}$、$2.40\times10^{-4}$ kg/kg（末行差异主要来自两套离散的 $t^*$ 相差 0.12--0.26 h）；$0.125\to0.0625$ mm 的对齐行最大差为 $1.03\times10^{-4}$，0.125 mm 解的 Richardson 估计（按 $t^*$ 链实测阶 $p=1.22$）约 $7.8\times10^{-5}$，仍高于 $0.5\times10^{-4}$ 阈值——与问题 2 不同，表 5 末行由 $t^*$ 时间插值取值、且最大差出现在 42 h、1.5 cm 处的陡梯度区，实测收敛阶 1.1--1.3、低于格式二阶，故需再加密一档。实测 $0.0625\to0.03125$ mm：对齐行最大差 $4.14\times10^{-5}$ kg/kg（最大差仍位于 42 h、1.5 cm 处），0.0625 mm 解的 Richardson 估计（按实测阶 $p=1.33$）约 $2.7\times10^{-5}$，低于阈值（表 5 逐格比对：0.0625 mm 解较 0.125 mm 解有 19 个单元格第四位变化 $+1\times10^{-4}$，印证 0.125 mm 解第四位不可靠）。$t^*$ 随步长的序列为 $56.7064/56.9672/57.0874/57.1392/57.1597$ h（$0.5/0.25/0.125/0.0625/0.03125$ mm），实测收敛阶 1.12/1.22/1.33，按末段实测阶外推 $57.1732$ h，与生产值相差 0.034 h，空间离散不确定度约 0.03 h。故取 $\Delta r=0.0625$ mm（$N=321$），$\Delta t=2.5$ s（减半后前 6 h 最大差 $5.58\times10^{-6}$）。

\textbf{问题 4。}$\zeta$ 坐标下同样按减半加密实测：$1/20\to1/40\to1/80$（$N=21,41,81$），表 6 的 6 h 行最大差依次 $5.33\times10^{-4}$、$1.62\times10^{-4}$ kg/kg（固定距离列线性插值取值）；线性插值在陡梯度区的 $O(\Delta\zeta^2)$ 误差主导了行差（实测 $1/80\to1/160$ 差 $1.58\times10^{-4}$、收敛比仅 2.5），故固定距离列改用 Pchip 三次插值（解光滑，误差 $O(\Delta\zeta^3)$），此后 $1/80\to1/160$ 最大差降至 $1.48\times10^{-4}$、$1/160\to1/320$ 为 $5.33\times10^{-5}$ kg/kg（实测收敛阶约 1.5），$1/160$ 离散的 Richardson 估计约 $3.0\times10^{-5}$，低于 $0.5\times10^{-4}$ 阈值（$1/80$ 离散约 $8.4\times10^{-5}$，不达标），故取 $\Delta\zeta=1/160$（$N=161$）；$t^*$ 序列 $50.5416/50.6928/50.7713/50.8049/50.8170$ h（$1/20\sim1/320$），按实测阶（末三段 $p=1.47$）外推 $50.8239$ h，最细两档差 0.0122 h，空间离散不确定度约 0.01 h。$\Delta t=5$ s（减半后最大差 $1.10\times10^{-5}$）。

以上加密过程均为实测数据：输出行在平滑区收敛比达理论值 4（问题 1、2 实测比值 4.0）、误差按二次规律下降，陡梯度区与插值取值处收敛阶偏低（问题 3、4，改用实测阶估计，7.2 节汇总），并与 7.1 节的解析解对照、7.3 节的守恒检验交叉印证。





\section{六、求解结果}

以下结果均按题面要求保留四位小数：问题 1、2 给出各时刻的温度与水分浓度两套数据；问题 3、4 给出各时刻的水分浓度，并附烘干结束时刻 $t^*$（205701 s / 182897 s）的一行，与表 5、表 6 的末行一致。

\subsection{6.1 问题 1：预热平衡阶段（表 1、表 2）}

空间步长 $\Delta r=0.03125$ mm、时间步长 $\Delta t=0.125$ s，附件 1 时变边界下求解至 1800 s，结果如表 1、表 2 所示。

\begin{center}
{\heiti\zihao{-4} 表 1\quad 30 分钟内药材的温度（$^\circ$C）}
\par\medskip
{\zihao{5}\begin{tabular}{cccccc}
\toprule
时间/s & \multicolumn{5}{c}{到药材中心的距离/cm} \\
\cmidrule(lr){2-6}
 & 0 & 0.5 & 1 & 1.5 & 2 \\
\midrule
100  & 28.0001 & 28.0003 & 28.0040 & 28.0327 & 28.1801 \\
300  & 28.0408 & 28.0635 & 28.1514 & 28.3681 & 28.8489 \\
600  & 28.4533 & 28.5360 & 28.8039 & 29.3159 & 30.1652 \\
900  & 29.3243 & 29.4583 & 29.8755 & 30.6161 & 31.7304 \\
1200 & 30.5427 & 30.7098 & 31.2223 & 32.1126 & 33.4276 \\
1500 & 31.9957 & 32.1867 & 32.7660 & 33.7463 & 35.1203 \\
1800 & 33.5753 & 33.7720 & 34.3642 & 35.3621 & 36.7856 \\
\bottomrule
\end{tabular}}
\end{center}

\newpage
\begin{center}
{\heiti\zihao{-4} 表 2\quad 30 分钟内药材的水分浓度（kg/kg）}
\par\medskip
{\zihao{5}\begin{tabular}{cccccc}
\toprule
时间/s & \multicolumn{5}{c}{到药材中心的距离/cm} \\
\cmidrule(lr){2-6}
 & 0 & 0.5 & 1 & 1.5 & 2 \\
\midrule
100  & 2.5500 & 2.5500 & 2.5500 & 2.5500 & 2.2471 \\
300  & 2.5500 & 2.5500 & 2.5500 & 2.5492 & 2.0509 \\
600  & 2.5500 & 2.5500 & 2.5500 & 2.5353 & 1.8770 \\
900  & 2.5500 & 2.5500 & 2.5497 & 2.5045 & 1.7547 \\
1200 & 2.5500 & 2.5500 & 2.5482 & 2.4646 & 1.6586 \\
1500 & 2.5500 & 2.5499 & 2.5445 & 2.4206 & 1.5787 \\
1800 & 2.5500 & 2.5497 & 2.5383 & 2.3755 & 1.5102 \\
\bottomrule
\end{tabular}}
\end{center}

结果分析：温度场自外向内升温，表面 1800 s 达 36.79 $^\circ$C，中心仅 33.58 $^\circ$C，径向温差 3.2 $^\circ$C；1800 s 时烘房温度 41.51 $^\circ$C，表面温度尚未达到烘房温度。

水分场的变化集中于近表面约 5 mm 的边界层内（特征长度 $\delta\sim D/h_m\approx5$ mm），表面水分浓度由 2.55 kg/kg 降至 1.51 kg/kg，而中心在 1800 s 内基本不变（2.5500 kg/kg）。这一「外层先失水、内部未及响应」的形态正是预热平衡阶段的特征——水分扩散系数仅约 $(4\sim5)\times10^{-9}$ m$^2$/s，水分扩散深度 $\sqrt{Dt}\approx2.6$ mm，远小于半径，故内部含水率尚未显著下降。

\subsection{6.2 问题 2：前 3 小时（表 3、表 4）}

按附录 3 经验式与附件 1 边界求解，结果如表 3、表 4 所示。

\begin{center}
{\heiti\zihao{-4} 表 3\quad 3 小时内药材的温度（$^\circ$C）}
\par\medskip
{\zihao{5}\begin{tabular}{cccccc}
\toprule
时间/h & \multicolumn{5}{c}{到药材中心的距离/cm} \\
\cmidrule(lr){2-6}
 & 0 & 0.5 & 1 & 1.5 & 2 \\
\midrule
0.5 & 32.1893 & 32.3821 & 32.9660 & 33.9608 & 35.4131 \\
1.0 & 40.3816 & 40.5536 & 41.0605 & 41.8782 & 42.9976 \\
1.5 & 45.8468 & 45.9348 & 46.1930 & 46.6051 & 47.1401 \\
2.0 & 48.4502 & 48.4881 & 48.5985 & 48.7736 & 49.0033 \\
2.5 & 49.4670 & 49.4792 & 49.5137 & 49.5653 & 49.6609 \\
3.0 & 49.8495 & 49.8553 & 49.8746 & 49.9101 & 49.9664 \\
\bottomrule
\end{tabular}}
\end{center}

\begin{center}
{\heiti\zihao{-4} 表 4\quad 3 小时内药材的水分浓度（kg/kg）}
\par\medskip
{\zihao{5}\begin{tabular}{cccccc}
\toprule
时间/h & \multicolumn{5}{c}{到药材中心的距离/cm} \\
\cmidrule(lr){2-6}
 & 0 & 0.5 & 1 & 1.5 & 2 \\
\midrule
0.5 & 2.5499 & 2.5489 & 2.5255 & 2.3257 & 1.6486 \\
1.0 & 2.5257 & 2.4948 & 2.3578 & 2.0230 & 1.4711 \\
1.5 & 2.3860 & 2.3256 & 2.1344 & 1.8020 & 1.3476 \\
2.0 & 2.1708 & 2.1084 & 1.9236 & 1.6259 & 1.2311 \\
2.5 & 1.9566 & 1.9006 & 1.7360 & 1.4720 & 1.1166 \\
3.0 & 1.7662 & 1.7165 & 1.5702 & 1.3333 & 1.0081 \\
\bottomrule
\end{tabular}}
\end{center}

结果分析：温度场约 2 h 起趋于均匀并逼近烘房温度（3 h 时中心 49.85 $^\circ$C、表面 49.97 $^\circ$C），传热时间尺度（热扩散率 $\alpha=k/(\rho c_p)\approx1.5\times10^{-7}$ m$^2$/s，$R^2/\alpha\approx2800$ s）远小于干燥时间尺度，说明整个烘干过程中药材温度在数小时内即与环境平衡，此后的干燥是典型的「等温干燥」过程。

水分场方面，3 h 时表面降至 1.0081 kg/kg、中心降至 1.7662 kg/kg，径向含水率梯度达 0.76 kg/kg，干燥前沿已明显向内部推进。对比问题 1：附录 3 的 $D$ 比附录 2 的 $D$ 大约 2--3 倍（干燥前期典型含水率下；低含水率端差距更大），且随升温进一步增大，故 3 h 内的失水幅度显著大于预热段。

\subsection{6.3 问题 3：烘干时长与全程水分场（表 5）}

以 $\max_r C(r,t)\le0.15$ kg/kg 为判据求解至全程，得
\begin{equation}
t^*=205701\ \mathrm{s}=\mathbf{57.14\pm0.03\ h}\ \ (\text{约 2.38 天；不确定度估计见 7.2 节}),
\end{equation}
与题述「烘干过程一般持续 2--3 天」一致。每隔 6 h 的水分浓度如表 5 所示。

\begin{center}
{\heiti\zihao{-4} 表 5\quad 药材烘干过程的水分浓度（kg/kg）}
\par\medskip
{\zihao{5}\begin{tabular}{cccccc}
\toprule
时间/h & \multicolumn{5}{c}{到药材中心的距离/cm} \\
\cmidrule(lr){2-6}
 & 0 & 0.5 & 1 & 1.5 & 2 \\
\midrule
6     & 1.0155 & 0.9868 & 0.9006 & 0.7546 & 0.5334 \\
12    & 0.4548 & 0.4418 & 0.4019 & 0.3289 & 0.1640 \\
18    & 0.2979 & 0.2906 & 0.2679 & 0.2247 & 0.0843 \\
24    & 0.2371 & 0.2321 & 0.2161 & 0.1851 & 0.0663 \\
30    & 0.2053 & 0.2013 & 0.1887 & 0.1637 & 0.0597 \\
36    & 0.1854 & 0.1820 & 0.1714 & 0.1500 & 0.0565 \\
42    & 0.1716 & 0.1687 & 0.1593 & 0.1403 & 0.0547 \\
48    & 0.1614 & 0.1588 & 0.1503 & 0.1331 & 0.0536 \\
54    & 0.1535 & 0.1511 & 0.1433 & 0.1274 & 0.0528 \\
烘干结束时间（57.14 h） & 0.1500 & 0.1477 & 0.1402 & 0.1249 & 0.0525 \\
\bottomrule
\end{tabular}}
\end{center}

结果分析：干燥过程呈明显的两阶段形态。前 12 h 含水率快速下降（表面由 2.55 降至 0.164 kg/kg），此阶段表面传质与内部扩散\textbf{共同控制}——传质 Biot 数 $Bi_m=h_mR/D(C_R,T_R)$ 在干燥初期约为 2.8（$C_R=2.55$ kg/kg、$T_R\approx28$ $^\circ$C）、6 h 时约 2.4，内部扩散阻力已为表面传质阻力的 2--3 倍；此后进入降速干燥段——表面 $C_R$ 降至 0.164 kg/kg 使 $D$ 按 $\mathrm{e}^{-0.45/C}$ 急剧衰减，$Bi_m$ 于 12 h 增至约 15、18 h 已超 200，内部扩散完全主导，表面含水率 54 h 时仅 0.0528 kg/kg，已逼近烘房水分浓度 0.04986 kg/kg（表面传质驱动力趋于消失）。

烘干判据由中心点决定（中心始终是含水率最高处），最终 $t^*=57.14$ h。

\subsection{6.4 问题 4：含收缩的烘干时长（表 6）}

附件 2 的 $R(t)$ 经 Pchip 插值代入 $\zeta$ 坐标离散求解，得
\begin{equation}
t^*=182897\ \mathrm{s}=\mathbf{50.80\pm0.01\ h}\ \ (\text{约 2.12 天；不确定度估计见 7.2 节}).
\end{equation}
每隔 6 h 的水分浓度如表 6 所示。

\begin{center}
{\heiti\zihao{-4} 表 6\quad 含收缩的药材烘干过程水分浓度（kg/kg）}
\par\medskip
{\zihao{5}\begin{tabular}{ccccccc}
\toprule
时间/h & \multicolumn{6}{c}{到药材中心的距离/cm} \\
\cmidrule(lr){2-7}
 & 0 & 0.5 & 1 & 1.5 & 2 & 药材表面 \\
\midrule
6     & 1.7164 & 1.5350 & 1.0212 & --- & --- & 0.4215 \\
12    & 0.7341 & 0.6516 & 0.4068 & --- & --- & 0.1666 \\
18    & 0.4066 & 0.3668 & 0.2387 & --- & --- & 0.0889 \\
24    & 0.2837 & 0.2598 & 0.1783 & --- & --- & 0.0671 \\
30    & 0.2253 & 0.2085 & 0.1487 & --- & --- & 0.0593 \\
36    & 0.1920 & 0.1789 & 0.1312 & --- & --- & 0.0558 \\
42    & 0.1706 & 0.1598 & 0.1196 & --- & --- & 0.0539 \\
48    & 0.1556 & 0.1463 & 0.1112 & --- & --- & 0.0529 \\
烘干结束时间（50.80 h） & 0.1500 & 0.1413 & 0.1080 & --- & --- & 0.0525 \\
\bottomrule
\end{tabular}}
\end{center}

结果分析：半径在干燥前期快速收缩（前 3.5 h 内由 2 cm 缩至约 1.5 cm），故 1.5 cm、2 cm 两列自 6 h 起即已落在药材之外。与问题 3 对比：$t^*$ 由 57.14 h 缩短为 50.80 h，缩短约 6.3 h（11.1\%）。这一净效应由两个相反方向的因素合成：附录 4 的扩散系数比附录 3 小约 2--5 倍（高含水率端约 5 倍、低含水率端约 2 倍），使干燥变慢；而收缩使算子系数 $D/R^2$ 增大 2.8 倍（等效于扩散路径缩短），使干燥加快。

由于问题 3 与问题 4 同时更换了物性与几何，两因素贡献须用同一套离散设置的消融对照分离，结果如表 7 所示。四个组合采用同一套验证离散（$\Delta\zeta=1/160$、物理分辨率 0.125 mm、$\Delta t=5$ s，迭代判据与 5.5 节相同）；附录 3 固定几何值 57.0884 h 与问题 3 生产值（57.14 h，0.0625 mm）相差 0.05 h，属该验证离散的空间偏差，各组合受同向影响，表内百分比结论的影响小于 0.1 个百分点。

\begin{center}
{\heiti\zihao{-4} 表 7\quad 物性 × 几何的 2×2 消融对照：烘干时长 $t^*$/h}
\par\medskip
{\zihao{5}\begin{tabular}{lccc}
\toprule
 & 固定半径 $R=2$ cm & 收缩 $R(t)$（附件 2） & 收缩净效应 \\
\midrule
附录 3 物性 & 57.0884 & 25.1162 & $-31.9722$（$-56.0\%$） \\
附录 4 物性 & 129.0164 & 50.8049 & $-78.2115$（$-60.6\%$） \\
\bottomrule
\end{tabular}}
\end{center}

表中可见：固定几何下附录 3→附录 4 的物性差异使烘干时长增加 71.93 h，增幅为 $+126.0\%$；收缩在附录 3、附录 4 下分别使时长缩短 31.97 h（$-56.0\%$）与 78.21 h（$-60.6\%$），且 $D$ 越小、干燥越受内部扩散控制时，收缩的加速作用越强；两因素叠加（问题 3→问题 4 的实际组合）的净效应为 $-6.3$ h（约 $-11\%$），收缩加速占主导。

此外，表面含水率列显示表面在 12 h 内即降至 0.167 kg/kg 并持续逼近烘房水分浓度 0.04986 kg/kg，失水主要发生在干燥前期，与收缩曲线（前期收缩快、67 h 起稳定于 1.198 cm）互为印证。

\section{七、模型的检验}

\subsection{7.1 Bessel 级数解析解对照（问题 1 热传导）}

对问题 1 的热传导方程 \eqref{eq:heat1}，在常数边界 $T_a=41.513\ ^\circ$C（附件 1 中 1800 s 时刻的环境温度）下存在解析解（Carslaw--Jaeger 级数解\cite{carslaw}）：
\begin{equation}
T(r,t)=T_a+(T_0-T_a)\sum_{n=1}^{\infty}A_n\,\mathrm{e}^{-\alpha\beta_n^2 t/R^2}J_0\!\Bigl(\frac{\beta_n r}{R}\Bigr),
\label{eq:bessel}
\end{equation}
其中 $\beta_n$ 为特征方程 $\beta J_1(\beta)=\mathrm{Bi}\,J_0(\beta)$ 的正根（$\mathrm{Bi}=hR/k=1.389$），系数 $A_n=\dfrac{2J_1(\beta_n)}{\beta_n\bigl(J_0^2(\beta_n)+J_1^2(\beta_n)\bigr)}$。取前 40 项级数与相同常数边界下的数值解（空间步长 0.03125 mm）在 $t=1800$ s 对照：中心点数值解 37.8050 $^\circ$C、解析解 37.8050 $^\circ$C；表面数值解 39.4483 $^\circ$C、解析解 39.4483 $^\circ$C，四位小数完全一致（误差不超过 $5\times10^{-5}\ ^\circ$C）。这同时验证了离散格式与 Robin 表面处理（含表面控制体体积、通量系数 $Rh$）的正确性。

\subsection{7.2 收敛性检验}

\textbf{空间收敛}。问题 1 在 $\Delta r=0.125,0.0625,0.03125$ mm 三种空间步长（$\Delta t=0.125$ s 固定、同一位置采样）下求表 1、表 2：表 1 最大差依次为 $2.30\times10^{-5}$、$5.75\times10^{-6}$ $^\circ$C，比值 4.0；表 2 依次为 $3.51\times10^{-4}$、$8.72\times10^{-5}$ kg/kg，比值 4.0，最大差位于 $(100\ \mathrm{s},$ 表面$)$ 的浓度边界层处。差值比接近理论值 4（二阶格式的理论收敛比，依据 5.5 节的截断误差分析），验证格式为空间二阶收敛；0.03125 mm 步长下表 2 的 Richardson 残余误差约 $2.9\times10^{-5}$，四位小数可信。问题 2 在 0.5/0.25/0.125 mm 三种空间步长下，表 4 最大差依次为 $1.86\times10^{-4}$、$4.66\times10^{-5}$ kg/kg，比值 4.0，同样二阶收敛（表 3 对应差值为 $2.08\times10^{-4}$、$5.21\times10^{-5}\ ^\circ$C）。

问题 3 在 0.125/0.0625/0.03125 mm 下的表 5 对齐行逐格差：$0.125\to0.0625$ mm 最大 $1.03\times10^{-4}$、$0.0625\to0.03125$ mm 最大 $4.14\times10^{-5}$ kg/kg，最大差均位于 42 h、1.5 cm 处；表 5 行的实测收敛阶低于格式二阶（末行由 $t^*$ 时间插值取值、最大差位置处于时间-浓度剖面最陡处），按末段实测阶 $p=1.33$ 估计，0.0625 mm 解的残余误差约 $2.7\times10^{-5}$，四位小数可信。问题 4 的 $\zeta$ 方向加密实测：固定距离列经线性插值取值时，$1/40\to1/80\to1/160$ 的表 6 行最大差为 $3.98\times10^{-4}$、$1.58\times10^{-4}$ kg/kg，收敛比仅 2.5——线性插值的 $O(\Delta\zeta^2)$ 误差在 30 h 前后的陡梯度区主导了行差；固定距离列改用 Pchip 三次插值（$C(\zeta)$ 为抛物方程的光滑解，插值误差 $O(\Delta\zeta^3)$）后，$1/80\to1/160\to1/320$ 的最大差降至 $1.48\times10^{-4}$、$5.33\times10^{-5}$ kg/kg，实测收敛阶约 1.5，$1/160$ 离散的 Richardson 估计约 $3.0\times10^{-5}$，四位小数可信。

烘干时长 $t^*$ 的空间离散敏感性：问题 3 在 $\Delta r=0.5$、0.25、0.125、0.0625、0.03125 mm 五档下 $t^*$ 为 56.7064、56.9672、57.0874、57.1392、57.1597 h，实测收敛阶 1.12/1.22/1.33，按末段实测阶外推 57.1732 h，与生产值相差 0.034 h，空间离散不确定度约 0.03 h；问题 4 在 $\Delta\zeta=1/20\sim1/320$ 五档下 $t^*$ 为 50.5416、50.6928、50.7713、50.8049、50.8170 h，按实测阶 $p=1.47$ 外推 50.8239 h，不确定度约 0.01 h。

\textbf{时间收敛}。问题 1 时间步加密（$\Delta t$ 减半）后表 2 差值为 $O(10^{-5})$；问题 2：$\Delta t=0.5\to0.25$ s，表 3 最大差 $2.02\times10^{-5}$ $^\circ$C、表 4 最大差 $6.59\times10^{-6}$ kg/kg；问题 3：$\Delta t=2.5\to1.25$ s 前 6 h 最大差 $5.58\times10^{-6}$ kg/kg；问题 4：$\Delta t=5\to2.5$ s 最大差 $1.10\times10^{-5}$ kg/kg。时间离散误差均比空间误差小 1--2 个量级，不影响四位小数。

\subsection{7.3 离散质量守恒检验}

按恒等式 \eqref{eq:conserv} 检验全程离散守恒（问题 1 取 1800 s，问题 3、4 取至 $t^*$）：问题 1 在 0.03125 mm 步长下 $\sum v_i\Delta C_i=-5.128889\times10^{-5}$，与 $-R\int h_m(C_R-C_a)\dd t$ 完全相等，闭合差 $3.9\times10^{-19}$；问题 3 闭合差 $6.31\times10^{-11}$（$\sum v_i\Delta C_i=-4.851549\times10^{-4}$；末行由 $t^*$ 时间插值取值，与离散时间层不完全配对，闭合差相对值 $1.3\times10^{-7}$）；问题 4（动边界，按物理体积加权求和、体积收缩引入几何项 $2\dot{R}C/R$）闭合差 $4.1\times10^{-7}$。问题 1 达机器精度、问题 3 达 $10^{-11}$ 量级，验证了「表面通量作用于真实表面节点」离散方案的离散通量严格闭合性；问题 4 的残差源于动边界下 $R$ 取半步中点与梯形积分配对的 $O(\Delta t^2)$ 配对误差，相对总失水量仅 $8\times10^{-4}$，可忽略。

守恒量按 $\sum v_iC_i$ 计，即水分浓度对体积的积分——它是 $C$-方程本身的离散不变量，本检验针对数值格式的守恒性；物理总水量为 $\int\rho_d C\,\dd V$（$\rho_d=\rho/(1+C)$ 为干物质密度），其中 $\rho_d$ 随 $C$ 的变化已由题给有效扩散系数的标定隐含吸收（假设 H5），故物理水量守恒不能由 $\sum v_iC_i$ 直接检验，本检验的结论限于离散格式层面。

\subsection{7.4 渐近行为检验}

以常数边界 $(T_a,C_a)=(50.165\ ^\circ\mathrm{C},0.04986)$ 长时程求解：问题 2/3/4 的模型在 $t=10^5$ s 时中心温度即达 50.1650 $^\circ$C，与 $T_a$ 四位小数一致；表面水分浓度由 0.0616（$10^5$ s）缓慢趋近 0.04986（$3\times10^5$ s 时为 0.0512），渐近速率随 $D\to0$ 而衰减，与平衡态 $C\to C_a$ 的物理预期一致（$C=C_a$ 时式 \eqref{eq:D3}、\eqref{eq:D4} 的驱动力与扩散系数同时趋零，故呈指数拖尾）。问题 1 模型（$T_a=41.513\ ^\circ$C）同样精确收敛：中心温度 $2\times10^4$ s 起恒为 41.5130 $^\circ$C。

\subsection{7.5 参数物理合理性检验}

题给扩散系数中的 Arrhenius 指数 $-3850/T$ 对应表观活化能
\begin{equation}
E_a=3850\,R_g=3850\times8.3145\ \mathrm{J/mol}\approx32.0\ \mathrm{kJ/mol}.
\end{equation}
文献报道的中药材热风干燥活化能：丹参 $D_{\mathrm{eff}}=1.7\times10^{-8}\sim9.4\times10^{-8}$ m$^2$/s、$E_a=40.0$ kJ/mol\cite{danshen}；鲜地黄 $E_a=34.94$ kJ/mol\cite{dihuang}；白术 $E_a=37.47$ kJ/mol\cite{baizhu}；题给 $D$ 的数值范围 $10^{-11}\sim10^{-8}$ m$^2$/s 亦与上述实测同量级。国内对中药干燥过程的传热传质数值模拟亦有报道\cite{wang}。$E_a=32.0$ kJ/mol 高于纯液态水自扩散活化能（18--20 kJ/mol\cite{mills}）而低于化学吸附能级，处于解吸-汽化控制的合理区间，说明题给经验式在物理上自洽。此外 $t^*=57.14$ h、50.80 h（约 2.1--2.4 天）均落在题述「2--3 天」范围内，与题面经验时长在数量级上一致；题面时长并非独立实验数据，此处仅作数量级合理性参照，不构成独立验证。

\subsection{7.6 汽化潜热影响的数量级分析}

水分汽化在表面吸热，可能使表面温度低于环境温度。按题设的有效模型观点（假设 H5），主模型暂不显式引入潜热汇项；为定量说明显式计入潜热可能带来的影响，作如下数量级估计。

\textbf{(1) 湿球温度估计。}按主模型口径，$C_a$ 取题给单位 kg/kg、与干基含水率 $C$ 同纲，解释为烘房空气状态对应的药材等效平衡水分浓度（见 5.1 节），故 Robin 边界驱动力 $C_R-C_a$ 直接成立，$C\to C_a$ 即干燥终点的平衡态。为估计显式计入潜热时表面温度的下限，本节另作\textbf{情景假设}：若将 $C_a$ 按空气含湿比解释，则 50 $^\circ$C 下空气的饱和含湿量约 0.0868 kg/kg\cite{incropera}，相对湿度 $\mathrm{RH}\approx57\%$，由焓湿图求得湿球温度 $T_{wb}\approx41.5\ ^\circ$C；该值仅作为情景下限，两种解释下本节的数量级结论一致。

\textbf{(2) 蒸发吸热与对流显热的比值。}单位面积上蒸发带走的热流密度与对流供热的显热流密度分别为
\begin{equation}
q_{\mathrm{evap}}=\rho_d\,\Delta H_v\,h_m\,(C_R-C_a),\qquad q_{\mathrm{sens}}=h\,(T_a-T_R),
\label{eq:qratio}
\end{equation}
其中 $\rho_d=\rho(C)/(1+C)$ 为干物质密度（$C$ 以干基计）。取 $\Delta H_v(50\,^\circ\mathrm{C})\approx2.38\times10^6$ J/kg\cite{incropera}、$h_m=8\times10^{-7}$ m/s、$h=25$ W/(m$^2\!\cdot$K)，在干燥中期典型值 $C_R-C_a\approx1$ kg/kg、$T_a-T_R\approx10$ K 下，附录 3 的干物质密度在干燥中期（$C\approx0.5\sim1$）为 $\rho_d\approx400\sim500$ kg/m$^3$、上限约 650 kg/m$^3$（$C\to0$ 时），代入：
\begin{equation}
\frac{q_{\mathrm{evap}}}{q_{\mathrm{sens}}}
=\frac{\rho_d\Delta H_v h_m(C_R-C_a)}{h(T_a-T_R)}
\approx\frac{(400\sim650)\times2.38\times10^6\times8\times10^{-7}\times1}{25\times10}\approx3\sim5.
\end{equation}
即蒸发吸热与对流显热同量级（比值约 3--5），若显式计入潜热汇，表面温度将明显低于环境、向湿球温度（约 41.5 $^\circ$C）靠拢。

\textbf{(3) 对干燥时长的影响。}若表面温度由 50 $^\circ$C 降至湿球温度约 41.5 $^\circ$C，Arrhenius 因子减小 $\mathrm{e}^{3850(1/323.3-1/314.7)}\approx0.72$，即扩散系数下降约 28\%，而降速干燥段 $t^*\propto R^2/D$，干燥时长约增加 $1/0.72-1\approx39\%$（$t^*$ 由 57.14 h 增至约 79.4 h，即约 3.3 天）；若表面进一步降至 40 $^\circ$C，则时长约增加 46\%（约 83.4 h，即约 3.5 天）。

可见若显式计入潜热，预测时长将超出题述「2--3 天」的范围——这从反面说明题给经验式不宜再显式叠加潜热项：$D(C,T)$、$h_m$ 等由真实干燥实验数据拟合，拟合得到的 $E_a=32.0$ kJ/mol 已含解吸与汽化的综合能垒（纯液态水扩散仅为 18--20 kJ/mol），潜热效应\textbf{可能}已按有效参数方式部分吸收，但吸收程度无法唯一判定；若再显式叠加 Luikov 型潜热汇项 $\rho_d\Delta H_v\,\partial C/\partial t$，存在与经验参数\textbf{重复计账}的风险，并将人为压低表面温度、高估干燥时长。故本文按题设采用有效模型，不显式引入潜热汇项，并将其影响以时长敏感性区间（约 $+40\%\sim50\%$）形式在结论中予以保留。

\subsection{7.7 参数灵敏度分析}

为评估模型结论对输入参数的稳健性，对问题 3、4 各取 8 个参数变体：烘房温度 $T_a\pm2\ ^\circ$C、烘房水分浓度 $C_a\pm20\%$、对流换热系数 $h\pm50\%$、对流传质系数 $h_m\pm50\%$，其余参数保持题给值，重算烘干时长，结果如表 8 所示（表中数值为各变体与基准烘干时长之差；该差值对空间离散不敏感，故在较粗的验证空间离散上计算：问题 3 取 0.25 mm、问题 4 取 $\Delta\zeta=1/80$）。

\begin{center}
{\heiti\zihao{-4} 表 8\quad 烘干时长 $t^*$ 的参数灵敏度}
\par\medskip
{\zihao{5}\begin{tabular}{lcc}
\toprule
参数变化 & 问题 3 $\Delta t^*$/h & 问题 4 $\Delta t^*$/h \\
\midrule
$T_a+2$ $^\circ$C & $-3.53$（$-6.2\%$） & $-3.14$（$-6.2\%$） \\
$T_a-2$ $^\circ$C & $+3.86$（$+6.8\%$） & $+3.43$（$+6.8\%$） \\
$C_a+20\%$ & $+0.36$（$+0.6\%$） & $+0.33$（$+0.6\%$） \\
$C_a-20\%$ & $-0.31$（$-0.5\%$） & $-0.24$（$-0.5\%$） \\
$h+50\%$ & $-0.03$（$-0.0\%$） & $-0.02$（$-0.0\%$） \\
$h-50\%$ & $+0.08$（$+0.1\%$） & $+0.07$（$+0.1\%$） \\
$h_m+50\%$ & $-2.04$（$-3.6\%$） & $-0.75$（$-1.5\%$） \\
$h_m-50\%$ & $+7.56$（$+13.3\%$） & $+3.34$（$+6.6\%$） \\
\bottomrule
\end{tabular}}
\end{center}

由表 8 可得如下结论：
\begin{enumerate}[label=(\arabic*),leftmargin=3em]
\item $t^*$ 对烘房温度最敏感，$T_a\pm2\ ^\circ$C 对应约 $\mp6\%\sim7\%$——温度通过 Arrhenius 因子直接提高扩散系数，恒温段温度控制是工艺优化的核心参数；
\item $h_m$ 减半在问题 3 中使 $t^*$ 增加 13.3\%，而问题 4 仅 6.6\%——附录 4 的 $D$ 更小，问题 4 的干燥更受内部扩散控制、表面传质系数的权重下降，与 6.4 节消融分析（物性差异 $+126\%$、收缩 $-60.6\%$）的机制一致；
\item $C_a$ 与 $h$ 的影响不足 1\%，烘房湿度波动与对流换热强度的不确定性对结论影响甚微；
\item 最不利联合工况（$T_a-2\ ^\circ$C 与 $h_m-50\%$ 同时施加，与表 8 同一验证离散）实际重算得问题 3 的 $t^*$ 增加 11.25 h（$+19.8\%$，68.22 h）、问题 4 增加 6.65 h（$+13.1\%$，57.42 h），与两因素单独影响之和 11.42 h、6.77 h 接近，非线性交互作用小；68.22 h 约合 2.84 天，仍在题述「2--3 天」范围内，模型结论稳健。
\end{enumerate}

\section{八、模型的评价与推广}

\subsection{8.1 模型优点}
\begin{enumerate}[label=(\arabic*),leftmargin=3em]
\item 控制方程严格取自题给条件：问题 1--3 用附录 2/3 参数、问题 4 用附录 4 参数与附件 2 收缩数据，四个问题的条件差异在 2.4 节逐一对照，各问参数来源严格对应。
\item 数值格式离散通量严格闭合：表面 Robin 通量直接作用于真实表面节点，与内部界面通量逐层对消，闭合差在 $10^{-7}$ 以内（问题 1 达机器精度 $10^{-18}$）。
\item 验证体系完整：解析解对照（Bessel 级数）、空间二阶收敛、守恒闭合、渐近行为、活化能文献互证五重检验全部通过。
\item 问题 4 的随体坐标处理使收缩模型无对流项，数值实现简洁且仍保持守恒性。
\end{enumerate}

\subsection{8.2 模型不足与改进方向}
\begin{enumerate}[label=(\arabic*),leftmargin=3em]
\item 未显式计入汽化潜热（见 7.6 节）；如需更精细的温湿耦合，可在获得题给参数之外的饱和蒸气压、吸湿等温线等数据后建立 Luikov 型模型\cite{whitaker,garcia}，与本文结果对照。
\item 径向一维模型忽略轴向与端面效应；对长径比更小的药材需建立二维轴对称模型\cite{mohan}。
\item 对流系数 $h,h_m$ 取常数；实际烘房风速、装载密度变化时可用经验关联式（如 Chilton--Colburn 类比\cite{incropera,defraeye}）修正。
\item 收缩假设各向同性且水分场径向一维；实际收缩可能呈各向异性并与局部含水率非线性相关。
\end{enumerate}

\subsection{8.3 结论与推广}

本文在题给有效参数体系下对四个问题给出定量结论：预热阶段（问题 1）1800 s 时表面温度 36.79 $^\circ$C、表面含水率 1.51 kg/kg，水分变化集中于近表面约 5 mm 的边界层；前 3 h（问题 2）温度场约 2 h 即趋于均匀并逼近烘房温度，此后干燥呈「等温干燥」形态；烘干时长（问题 3）$t^*=57.14\pm0.03$ h，含收缩时（问题 4）$t^*=50.80\pm0.01$ h。物性×几何消融表明收缩是两问差异的主导因素：收缩单独使时长缩短约 56\%--61\%，附录 4 物性使时长增加 126\%，净效应缩短约 11\%。工艺启示：烘房温度是最敏感参数（$T_a\pm2$ $^\circ$C 对应约 $\mp6\%\sim7\%$），传质系数次之，湿度与对流换热强度的影响不足 1\%，恒温段温度控制是干燥工艺优化的核心。

本文模型可直接推广为中药材干燥工艺的仿真工具：给定烘房温湿曲线与药材物性，即可预测任意时刻的温度场、水分场与烘干时长，进而为工艺参数（烘房温度、湿度、风速）的选取与能耗优化提供依据；模型输出的温度、水分时空分布可作为品质分析的定量输入——文献表明活性成分降解与干燥过程的温度、含水率直接相关\cite{ren}，故模型可用于评估变温策略（如分段降温保成分）中各时段的品质风险，但不直接预测成分含量。

\section*{附录：AI 工具使用声明}

\noindent 本参赛队在竞赛过程中使用了AI工具（Deepseek、Doubao），主要用于辅助题目分析与建模思路参考、代码调试和文字润色，所有结果均经本队人工核验；具体使用情况见AI工具使用详情说明。




\section*{附录：支撑材料文件列表}

\noindent 本文全部计算与结果可复现，支撑材料文件清单如下：
\begin{itemize}[leftmargin=2em]
\item result1.xlsx--result4.xlsx：四个问题的完整输出（问题 1、2 各含「温度」「水分浓度」两个工作表，问题 3、4 单表存水分浓度并含烘干结束时刻一行，均保留四位小数）；
\item q1\_solver.py、q2\_solver.py、q3\_solver.py、q4\_solver.py：四个问题的求解程序；
\item q78\_sensitivity.py：表 7 物性×几何消融与表 8 参数灵敏度的计算程序；
\item verify\_q1.py、verify\_q3.py、verify\_q4.py：Bessel 解析解对照、离散守恒闭合与收敛链的验证程序。
\end{itemize}

\begin{thebibliography}{9}
\setlength{\itemsep}{0pt}
\bibitem{crank} CRANK J. The Mathematics of Diffusion[M]. 2nd ed. Oxford: Clarendon Press, 1975.
\bibitem{carslaw} CARSLAW H S, JAEGER J C. Conduction of Heat in Solids[M]. 2nd ed. Oxford: Oxford University Press, 1959.
\bibitem{incropera} INCROPERA F P, DEWITT D P, BERGMAN T L, et al. Fundamentals of Heat and Mass Transfer[M]. 6th ed. New York: John Wiley \& Sons, 2006.
\bibitem{whitaker} WHITAKER S. Simultaneous heat, mass, and momentum transfer in porous media: a theory of drying[J]. Advances in Heat Transfer, 1977, 13: 119-203.
\bibitem{mujumdar} MUJUMDAR A S. Handbook of Industrial Drying[M]. 4th ed. Boca Raton: CRC Press, 2015.
\bibitem{simal} SIMAL S, ROSSELL\'O C, BERNA A, et al. Drying of shrinking cylinder-shaped bodies[J]. Journal of Food Engineering, 1998, 37(4): 423-435.
\bibitem{abbasi} ABBASI SOURAKI B, MOWLA D. Axial and radial moisture diffusivity in cylindrical fresh green beans in a fluidized bed dryer with energy carrier: modeling with and without shrinkage[J]. Journal of Food Engineering, 2008, 88(1): 9-19.
\bibitem{ren} 任迪峰, 毛志怀, 王建中. 中草药干燥过程中质量退化动力学模型的研究[J]. Chinese Journal of Chemical Engineering, 2004, 12(6): 822-825.
\bibitem{mohan} CHANDRA MOHAN V P, TALUKDAR P. Three dimensional numerical modeling of simultaneous heat and moisture transfer in a moist object subjected to convective drying[J]. International Journal of Heat and Mass Transfer, 2010, 53(21-22): 4638-4650.
\bibitem{defraeye} DEFRAEYE T, BLOCKEN B, CARMELIET J. Analysis of convective heat and mass transfer coefficients for convective drying of a porous flat plate by conjugate modelling[J]. International Journal of Heat and Mass Transfer, 2012, 55(1-3): 112-124.
\bibitem{garcia} GARC\'IA DEL VALLE J, SIERRA PALLARES J. Analytical solution for the coupled heat and mass transfer formulation of one-dimensional drying kinetics[J]. Journal of Food Engineering, 2018, 230: 99-113.
\bibitem{yang} 杨世铭, 陶文铨. 传热学[M]. 4 版. 北京: 高等教育出版社, 2006.
\bibitem{wang} 王学成, 康超超, 伍振峰, 等. 二至丸热风干燥过程温度均匀性模拟与实验[J]. 中草药, 2020, 51(5): 1226-1232.
\bibitem{danshen} LI Y Q, SHA X X, ZHANG Z, et al. Drying kinetics of Salviae Miltiorrhizae Radix et Rhizoma and dynamics of active components in drying process[J]. China Journal of Chinese Materia Medica, 2025, 50(1): 128-139.
\bibitem{dihuang} 孔兵, 郑琇梅, 栗嘉淇, 等. 鲜地黄干燥动力学及其对环烯醚萜类成分影响研究[J]. 人参研究, 2024, 36(6): 20-25.
\bibitem{baizhu} 郭慧玲, 徐梦甜, 伍振峰, 等. 白术热风干燥动力学及其挥发性成分变化规律的研究[J]. 中国中药杂志, 2022, 47(4): 922-930.
\bibitem{mills} MILLS R. Self-diffusion in normal and heavy water in the range 1-45$^\circ$C[J]. Journal of Physical Chemistry, 1973, 77(5): 685-688.
\end{thebibliography}

\end{document}

